\section*{第 \ref{ch:g8:speed-of-light} 章——光速；宇宙中的距离}

\begin{solution}{exo:g8:speed-of-light:1}
$c = \qty{300000}{km/s} = 3 \times 10^{8}\ \unit{m/s}$。山顶灯火
测得的只是人的反应：光穿过两山之间
所用的时间仅为百万分之一秒，远低于任何
人工计时的精度。
\end{solution}

\begin{solution}{exo:g8:speed-of-light:2}
木星的卫星按准点时刻发生食，是天上的
时钟。每次食的\emph{消息}只要地球
远离木星就晚到，靠近时则早到。
这一延迟测得的是光穿越地球\omterm{def:g6:sun-earth-moon:axisorbit}{轨道}
直径所用的时间，进而得出光速。
\end{solution}

\begin{solution}{exo:g8:speed-of-light:3}
\omterm{def:g2:moon-in-the-sky:moon}{月球}：$384000 \div 300000 \approx \qty{1.3}{s}$。太阳：
$150000000 \div 300000 = \qty{500}{s}$，即八分
二十秒。阳光总是展示太阳八分二十秒前的样子：
按其构造，这已是旧闻。
\end{solution}

\begin{solution}{exo:g8:speed-of-light:4}
一个距离：光在一年内所走过的路程，约
$9.5 \times 10^{12}$ 千米。
\end{solution}

\begin{solution}{exo:g8:speed-of-light:5}
四年。答案各异——四年前的你
还处在本书所涵盖的较低几个年级。
\end{solution}

\begin{solution}{exo:g8:speed-of-light:6}
单程约十一分钟左右，一次转向修正所回应的
画面至少已过时二十二分钟——漫游车会在
操纵杆的一次抖动抵达前就已陷入裂缝。因此
靠的是预先规划，而不是遥控操纵。
\end{solution}

\begin{solution}{exo:g8:speed-of-light:7}
\omterm{def:g2:moon-in-the-sky:moon}{月球}：约 \qty{1.3}{s}。太阳：八分钟。最近的\omterm{def:g4:solar-system:star}{恒星}：
四年。邻近星系：250 万年。
\end{solution}

\begin{solution}{exo:g8:speed-of-light:8}
光。\omterm{def:g7:motion-average-speed:formula}{速度} $c$ 被证明比任何金属杆或
国家标准尺都更稳定——因此尺子如今是由光导出的，
而不是用尺子去\omterm{def:g6:measurement-in-science:measuring}{测量}光。
\end{solution}

\begin{solution}{exo:g8:speed-of-light:9}
往返：$300000 \times 2.6 = 780000$ 千米；\omterm{def:g2:moon-in-the-sky:moon}{月球}处于
其一半处——\qty{390000}{km}。声呐的规律整体照用：
每次\omterm{ex:g4:how-sound-travels:echo}{回声}都付出了往返的路程，
取用之前先减半。
\end{solution}

\begin{solution}{exo:g8:speed-of-light:10}
灯：$3 \div (3 \times 10^{8}) = 10^{-8}$ 秒——
即一亿分之一秒。房间本身也是一份档案，
浅得难以察觉。接着是太阳（八分二十秒）、\omterm{def:g4:solar-system:star}{恒星}（一个
世纪）、星系（250 万年）——你望得越深，
页面就越古老。
\end{solution}

\begin{solution}{exo:g8:speed-of-light:11}
“上周二，一颗\omterm{def:g4:solar-system:star}{恒星}爆炸的光\emph{抵达了
我们}——爆炸本身发生在约 $8000$ 年前，
而它的消息从那时起便一直在路上。”
\end{solution}

\begin{solution}{exo:g8:speed-of-light:12}
$3 \times 10^{8}\ \text{千米} \div 1000\ \text{秒} = 3 \times
10^{5}$ \unit{km/s}——恰与现代的精确值一致（
历史上的估计，所用\omterm{def:g6:sun-earth-moon:axisorbit}{轨道}数据较粗略，
误差约在其四分之一以内：对 1600 年代而言，一个
令人惊叹的首个答案）。
\end{solution}

\begin{solution}{pb:g8:speed-of-light:1}
\textbf{1.} 你开口后约 \qty{1.3}{s}；最快的回应
要等往返全程，约
\qty{2.6}{s}。
\textbf{2.} 单程为 $36000 \div 300000 = \qty{0.12}{s}$；一问
一答要往返两次——接近半秒的
迟疑，恰好处于人耳的阈值。
\textbf{3.} “耀斑在屏幕亮起前八分二十秒就已爆发——
我们看到的是它\emph{曾经}的模样。”
\textbf{4.} \omterm{def:g2:moon-in-the-sky:moon}{月球}链路：\qty{1.3}{s} 对比 \qty{0.12}{s}
——大约是允许值的十倍。
\textbf{5.} $1.8 \times 10^{8} \div (3 \times 10^{5}) = 600$
秒：单程十分钟。
\textbf{6.} 往返共二十分钟：救援
画面已有十分钟之久，救援指令
也晚了十分钟——裂缝以十九分钟的优势胜出。漫游车
必须自带反应能力：车载紧急制动。
\textbf{7.} 最后一段须在 21:00 前到达；它在地球发出的时刻为
$21{:}00 - 10$ 分钟 $= 20{:}50$；传输
须提前三分钟开始，即 20:47 启动。
\textbf{8.} 延迟加倍，单程约
二十一分钟：\omterm{def:g4:solar-system:planet}{行星}并没有唯一确定的距离——两颗\omterm{def:g4:solar-system:planet}{行星}都在公转，
每一次对话都要针对一个移动的目标来安排。
\textbf{9.} $300000 \times 86400 \approx 2.6 \times 10^{10}$
千米——即二百六十亿千米；一问
一答要整整两天。
\textbf{10.} 光约四年抵达该\omterm{def:g4:solar-system:star}{恒星}；最快的
回复约在今晚的
发送之后八年才能到达。
\textbf{11.} “这束光离开它的星系，距今
已有 250 万年：这张快照播出的是地球
最早使用石器的时代——我们肉眼所能看到的最深
一页档案。”
\textbf{12.} 例如：“没有任何东西是瞬时的——
每个信号，包括我们自己的，最快也只以 $c$ 传播。
超出\omterm{def:g2:moon-in-the-sky:moon}{月球}之远，遥控操纵便敌不过延迟：远处的机器
必须自主思考。每台望远镜都是一位档案
阅读者——天体越远，消息越旧，
再也印不出更新的版本。”
\end{solution}
